\documentclass[10pt,leqno]{amsart}
\usepackage[utf8]{inputenc}
\usepackage[T1]{fontenc}
\usepackage[french]{babel}
\usepackage{amsmath,amssymb,amsthm}
\usepackage{mathrsfs}
\usepackage[all]{xy}
\usepackage{geometry}
\geometry{paperwidth=210mm,paperheight=297mm,textwidth=130mm,textheight=190mm,top=25mm,left=35mm}

\theoremstyle{plain}
\newtheorem{theorem}{Th\'eor\`eme}
\newtheorem{lemma}[theorem]{Lemme}
\newtheorem{proposition}[theorem]{Proposition}
\newtheorem{corollary}[theorem]{Corollaire}
\newtheorem*{itheorem}{Th\'eor\`eme}

\theoremstyle{definition}
\newtheorem{definition}[theorem]{D\'efinition}

\theoremstyle{remark}
\newtheorem{remark}[theorem]{Remarque}
\newtheorem*{conjecture}{Conjecture}

\DeclareMathOperator{\Pic}{Pic}
\DeclareMathOperator{\Spec}{Spec}
\DeclareMathOperator{\GL}{GL}
\DeclareMathOperator{\Gal}{Gal}
\DeclareMathOperator{\Coker}{Coker}
\DeclareMathOperator{\Hom}{Hom}
\DeclareMathOperator{\Aut}{Aut}
\DeclareMathOperator{\id}{id}
\DeclareMathOperator{\Tr}{Tr}
% \det already defined in amsmath
% \deg already defined in amsmath
\DeclareMathOperator{\pr}{pr}

\newcommand{\ZZ}{\mathbb{Z}}
\newcommand{\QQ}{\mathbb{Q}}
\newcommand{\CC}{\mathbb{C}}
\newcommand{\RR}{\mathbb{R}}
\newcommand{\PP}{\mathbb{P}}
\newcommand{\FF}{\mathbb{F}}
\newcommand{\GG}{\mathbb{G}}
\newcommand{\Oo}{\mathscr{O}}

\newcommand{\ua}{\underline{a}}

\begin{document}

\title{Les intersections compl\`etes de niveau de Hodge un}
\author{Pierre Deligne}
\address{(Bures-sur-Yvette)}
\thanks{\textit{Inventiones math.} \textbf{15}, 237--250 (1972) \copyright\ by Springer-Verlag 1972}

\maketitle

\section*{Introduction}
\addcontentsline{toc}{section}{Introduction}

Dans tout cet article, nous d\'esignerons par $n=2m+1$ un entier impair $>3$ fix\'e et par $\ua=(a_1,\ldots,a_\nu)$ une suite fix\'e d'entiers $>1$ tels que les intersections compl\`etes lisses $V_n(\ua)\subset\PP^{n+\nu}(\CC)$ de dimension $n$ et de multidegr\'e $\ua$ soient de niveau de Hodge $1$. Cette condition signifie que, pour une telle intersection compl\`ete $X$, intersection de $\nu$ hypersurfaces de degr\'es $a_1,\ldots,a_\nu$ qui, en les points de $X$, sont lisses et se coupent transversalement,
\begin{itemize}
\item[a)] pour $p+q=n$ et $|q-p|>1$, on a $H^q(X,\Omega^p)=0$;
\item[b)] $H^m(X,\Omega^{m+1})=0$.
\end{itemize}

Rappelons que, d'apr\`es Lefschetz,
\begin{itemize}
\item[c)] $H^*(X,\ZZ)$ est sans torsion;
\item[d)] pour $k$ impair, $k\neq n$, on a $H^k(X,\ZZ)=0$;
\item[e)] pour $p\le n$, $H^{2p}(X,\ZZ)$ est de rang $1$, purement de type $(p,p)$; si $\eta\in H^2(X,\ZZ)$ est la classe de cohomologie d'une section hyperplane, alors $H^{2p}(X,\QQ)=\QQ\cdot\eta^p$.
\end{itemize}

La liste des valeurs possibles de $(n,\ua)$ figure dans \cite{10}.

La jacobienne interm\'ediaire $J(X)$ de $X$ est par d\'efinition le tore complexe
\[
J(X)=H^n(X,\ZZ)\setminus H^n(X,\CC)/H^{m+1,m}(X)=H^{m,m+1}/H_n(X,\ZZ).
\]
Le cup produit sur $H^n(X,\ZZ)$ d\'efinit une forme de Riemann sur $J(X)$, qui est donc une vari\'et\'e ab\'elienne polaris\'ee de la s\'erie principale (=principalement polaris\'ee). On montre que la construction $X\mapsto J(X)$ est ``purement alg\'ebrique'', i.e. ind\'ependante de la topologie de $\CC$. En particulier, si $X$ est d\'efinie sur un corps $k\subset\CC$, alors $J(X)$ est d\'efinie sur le m\^eme corps $k$. Dans le cas o\`u $n=1$ (que nous n'excluons par la suite que pour des raisons de commodit\'e), $J(X)$ est la jacobienne de $X$. Pour les $V_3(3)$, $J(X)$ est la jacobienne de la surface de Fano de $X$ (voir \cite{1}). Pour $n=3$ et $X$ unirationnelle, les m\'ethodes de Manin \cite{8} permettent d'exprimer $J(X)$ \`a l'aide de jacobiennes. Pour les $V_n(2,2)$, si $(L_t)_{t\in\PP^1}$ est un pinceau de Lefschetz de sections hyperplanes, et si $C$ est le rev\^etement $n-1$ de $\PP^1$ de points les sous-espaces lin\'eaires de dimension $\frac{n-1}{2}$ des $L_t$, alors $J(X)$ est facteur direct \`a isog\'enie pr\`es de la jacobienne de $C$. Dans les autres cas, je ne connais pas d'interpr\'etation directe de $J(X)$. D'apr\`es Lieberman \cite{7}, il r\'esulte de la conjecture de Hodge que $J(X)$ est le quotient de l'ensemble des cycles alg\'ebriques sur $X$, de codimension $\frac{n+1}{2}$ et homologiquement \'equivalents \`a z\'ero, par une relation d'\'equivalence convenable.

Comme corollaire de la th\'eorie, on obtient (pour $n$ et $\ua$ comme fix\'es) que les intersections compl\`etes lisses de dimension $n$ et multidegr\'e $\ua$ sur un corps fini $\FF_q$ v\'erifient la conjecture de Weil. Pour les $V_5(3)$ et les $V_3(4)$, ce r\'esultat est, \`a ma connaissance, nouveau.

\section{Sur $\CC$}

\subsection{}\label{1.1}

Soit $X$ une vari\'et\'e alg\'ebrique complexe projective non singuli\`ere, connexe et de dimension $k$. La th\'eorie de dualit\'e des faisceaux alg\'ebriques coh\'erents fournit un isomorphisme (purement alg\'ebrique)
\begin{equation}\label{eq:1.1.1}
H^k(X,\Omega_X^k)\simeq\CC.
\end{equation}
Par ailleurs, la th\'eorie de Hodge (ici GAGA $+$ le lemme de Poincar\'e) fournit un isomorphisme
\begin{equation}\label{eq:1.1.2}
H^k(X,\Omega_X^k)\simeq H^{2k}(X,\CC).
\end{equation}
On sait que l'image par \eqref{eq:1.1.1} de la cohomologie enti\`ere $H^{2k}(X,\ZZ)\subset H^{2k}(X,\CC)$ est le r\'eseau $(2\pi i)^{-k}\ZZ\subset\CC$.

Ceci conduit \`a d\'efinir la structure de Hodge de Tate $\ZZ(s)$ comme la structure de Hodge de rang un, purement de type $(-s,-s)$, de r\'eseau entier $\ZZ(s)_\ZZ=(2\pi i)^s\ZZ\subset\ZZ(s)_\CC=\CC$. Le cup-produit d\'efinit un morphisme de structures de Hodge
\begin{equation}\label{eq:1.1.3}
H^k(X,\ZZ)\otimes H^k(X,\ZZ)\to\ZZ(-k).
\end{equation}
Les analogues de \eqref{eq:1.1.3} en cohomologie $l$-adique ou de De Rham sont les suivants.

\subsection{}\label{1.2}

Rappelons que $\ZZ/r(1)$ d\'esigne le groupe des racines $r$ i\`emes de l'unit\'e. On identifie $\ZZ(1)_\ZZ/r=_{\mathrm{dfn}}\ZZ(1)_\ZZ\otimes\ZZ/r\simeq\ZZ/r(1)$ par l'application $x\mapsto\exp(x/r)$. Pour tout nombre premier $\ell$, on pose
\[
\ZZ_\ell(1)=\varprojlim\ZZ/\ell^n(1)\qquad\text{(morphismes de transition $x\mapsto x^{\ell^{n-1}}$)},
\]
et $\ZZ_\ell(s)=_{\mathrm{dfn}}\ZZ_\ell(1)^{\otimes s}$.

L'application exponentielle identifie $\ZZ(1)_\ZZ\otimes\ZZ_\ell$ \`a $\ZZ_\ell(1)$ et $\ZZ(s)\otimes\ZZ_\ell$ \`a $\ZZ_\ell(s)$.

Si on tensorise \eqref{eq:1.1.3} par $\ZZ_\ell$, le morphisme
\begin{equation}\label{eq:1.2.1}
H^k(X,\ZZ_\ell)\otimes H^k(X,\ZZ_\ell)\to\ZZ_\ell(-k)
\end{equation}
obtenu n'est autre que le cup-produit en cohomologie $\ell$-adique. Il a un sens purement alg\'ebrique.

\subsection{}\label{1.3}

On a
\begin{equation}\label{eq:1.3.1}
H^k(X,\ZZ)\otimes\CC=H^k(X,\Omega^\bullet)=_{\mathrm{dfn}}H^k_{\mathrm{DR}}(X).
\end{equation}
Ici et par la suite, nous identifierons toujours $\ZZ(k)_\CC=\ZZ(k)_\ZZ\otimes\CC$ avec $\CC$ par l'isomorphisme ``de d\'efinition'' qui envoie le r\'eseau entier sur $(2\pi i)^k\ZZ$. Tensorisons \eqref{eq:1.1.3} avec $\CC$. Le morphisme obtenu
\begin{equation}\label{eq:1.3.2}
H^k_{\mathrm{DR}}(X)\otimes H^k_{\mathrm{DR}}(X)\to\CC
\end{equation}
n'est autre que le cup-produit en cohomologie de De Rham. Il a un sens purement alg\'ebrique.

\subsection{}\label{1.4}

Prenons pour $X$ une intersection compl\`ete $V_n(\ua)$. Tensorisons \eqref{eq:1.1.3} par $\ZZ(2m)$ et multiplions le morphisme obtenu par $(-1)^m$. On obtient
\begin{equation}\label{eq:1.4.1}
H^n(X)(m)\otimes H^n(X)(m)\to\ZZ(-1).
\end{equation}
Par hypoth\`ese, la structure de Hodge $H^n(X)(m)$ est de type $(0,1)+(1,0)$. La forme $\psi_0$ est une polarisation de $H^n(X)(m)$ et, d'apr\`es la dualit\'e de Poincar\'e pour $X$, la forme altern\'ee $\psi_\ZZ$ sur le r\'eseau entier $H^n(X)(m)_\ZZ$ est de discriminant un.

Il existe donc une (unique) vari\'et\'e ab\'elienne polaris\'ee de la s\'erie principale $J(X)$, munie d'un isomorphisme de structures de Hodge polaris\'ees
\[
\alpha_\ZZ\colon H^n(X)(m)\xrightarrow{\sim} H^1(J(X)).
\]
On appelle $J(X)$ la \emph{jacobienne interm\'ediaire} de $X$.

\subsection{}\label{1.5}

Soit $S$ un sch\'ema lisse de type fini sur $\CC$ et soit $f\colon X\to S$ un morphisme propre et lisse dont les fibres sont des intersections compl\`etes $V_n(\ua)$. Alors, $R^nf_*\ZZ(m)$ est une variation de structures de Hodge polaris\'ees sur $S$, de type $(0,1)+(1,0)$ (voir \cite{5} \S\,2). D'apr\`es un th\'eor\`eme de Borel (voir \cite{4} 5.1 et 5.2 pour l'\'enonc\'e) elle d\'efinit un sch\'ema ab\'elien polaris\'e de la s\'erie principale $u\colon J(X/S)\to S$, muni d'un isomorphisme de variations de structures de Hodge polaris\'ees
\[
\alpha\colon R^nf_*\ZZ(m)\xrightarrow{\sim} R^1u_*\ZZ.
\]
Cette construction exprime que $J(X)$ varie alg\'ebriquement avec $X$; il ne serait pas difficile de l'\'etendre au cas o\`u $S$ est un sch\'ema sur $\CC$ quelconque.

Les analogues en cohomologie $\ell$-adique ou de De Rham de la variation de structures de Hodge $R^nf_*\ZZ(m)$ et de $\alpha$ sont les suivants.

\subsection{}\label{1.6}

Par tensorisation avec $\ZZ_\ell$, $\alpha$ d\'efinit un isomorphisme de $\ZZ_\ell$-faisceaux
\begin{equation}\label{eq:1.6.1}
\alpha_\ell\colon R^nf_*\ZZ_\ell(m)\xrightarrow{\sim} R^1u_*\ZZ_\ell.
\end{equation}
Le cup-produit en cohomologie $\ell$-adique, multipli\'e par $(-1)^m$, est un morphisme de $\ZZ_\ell$-faisceaux
\begin{equation}\label{eq:1.6.2}
\psi_\ell\colon R^nf_*\ZZ_\ell(m)\otimes R^nf_*\ZZ_\ell(m)\to\ZZ_\ell(-1)
\end{equation}
et $\alpha_\ell$ transforme $\psi_\ell$ en la forme de polarisation
\begin{equation}\label{eq:1.6.3}
\varphi_\ell\colon R^1u_*\ZZ_\ell\otimes R^1u_*\ZZ_\ell\to\ZZ_\ell(-1).
\end{equation}
Les formes \eqref{eq:1.6.2} et \eqref{eq:1.6.3} ont un sens purement alg\'ebrique.

\subsection{}\label{1.7}

Le faisceau analytique coh\'erent localement libre $R^nf_*\ZZ\otimes\Oo$ est sous-jacent au faisceau alg\'ebrique coh\'erent
\[
\mathscr{H}^n_{\mathrm{DR}}(X/S)=_{\mathrm{dfn}}R^nf_*\Omega^\bullet_{X/S}.
\]
La connexion int\'egrable \'evidente de $R^nf_*\ZZ\otimes\Oo$ s'identifie \`a la connexion de Gauss-Manin de $\mathscr{H}^n_{\mathrm{DR}}(X/S)$. Par tensorisation avec $\Oo$, $\alpha$ d\'efinit un isomorphisme de faisceaux analytiques
\begin{equation}\label{eq:1.7.1}
\alpha_{\mathrm{DR}}\colon\mathscr{H}^n_{\mathrm{DR}}(X/S)\xrightarrow{\sim}\mathscr{H}^1_{\mathrm{DR}}(J(X/S)/S).
\end{equation}
Les connexions de Gauss-Manin de $\mathscr{H}^n_{\mathrm{DR}}(X/S)$ et de $\mathscr{H}^1_{\mathrm{DR}}(J(X/S)/S)$ sont r\'eguli\`eres (\cite{5} 14 ou \cite{2} II \S7) et $\alpha_{\mathrm{DR}}$ est un isomorphisme horizontal. $\alpha_{\mathrm{DR}}$ est donc alg\'ebrique (\cite{2} II 5.9).

Le cup-produit en cohomologie de De Rham multipli\'e par $(-1)^m$, est un morphisme horizontal
\begin{equation}\label{eq:1.7.2}
\psi_{\mathrm{DR}}\colon\mathscr{H}^n_{\mathrm{DR}}(X/S)\otimes\mathscr{H}^n_{\mathrm{DR}}(X/S)\to\Oo_S
\end{equation}
et $\alpha_{\mathrm{DR}}$ transforme $\psi_{\mathrm{DR}}$ en la forme de polarisation
\begin{equation}\label{eq:1.7.3}
\varphi_{\mathrm{DR}}\colon\mathscr{H}^1_{\mathrm{DR}}(J(X/S)/S)\otimes\mathscr{H}^1_{\mathrm{DR}}(J(X/S)/S)\to\Oo_S.
\end{equation}

\subsection{}\label{1.8}

Soit $S_\ZZ$ le sous-sch\'ema du sch\'ema de Hilbert (sur $\ZZ$) qui param\'etrise les intersections compl\`etes lisses de dimension $n$ et multidegr\'e $\ua$ dans $\PP^{n+\nu}$. Soit $f\colon X_\ZZ\to S_\ZZ$, la famille ``universelle'' de $V_n(\ua)$ param\'etr\'ee par $S_\ZZ$. Pour tout anneau $A$, nous d\'esignerons par $f\colon X_A\to S_A$ le morphisme d\'eduit de $f$ par extension des scalaires de $\ZZ$ \`a $A$.

Posons $G=PGL_{n+\nu+1}$. Le sch\'ema en groupe $G$ (sur $\ZZ$) agit sur $X_\ZZ/S_\ZZ$ via son action sur $\PP^{n+\nu}$.

\begin{proposition}\label{prop:1.9}
Soit $g\colon Y\to T$ un morphisme propre et lisse dont les fibres g\'eom\'etriques sont des intersections compl\`etes $V_n(\ua)$.
\begin{enumerate}
\item[(i)] Localement sur $T$ pour la topologie \'etale, $Y$ est isomorphe \`a l'intersection de $\nu$ hypersurfaces $H_1\cap\ldots\cap H_\nu$ de degr\'es $a_1,\ldots,a_\nu$ dans l'espace projectif $\PP^{n+\nu}$ sur $T$.
\item[(ii)] En particulier, localement sur $T$ pour la topologie \'etale, il existe $u\colon T\to S_\ZZ$ et un $T$-isomorphisme $t\colon Y\xrightarrow{\sim}u^*(X_\ZZ/S_\ZZ)$.
\item[(iii)] Si $(u',t')$ est un autre syst\`eme comme en \textup{(ii)}, alors, localement sur $T$ pour la topologie \'etale, il existe $a\in G(T)$ tel que $(u',t')=a\cdot(u,t)$.
\item[(iv)] $S_\ZZ$ est lisse sur $\Spec(\ZZ)$.
\end{enumerate}
\end{proposition}

Dans ce paragraphe, le seul corollaire de \ref{prop:1.9} que nous n'utiliserons est la lissit\'e de $S_\CC$. On pourrait m\^eme s'en dispenser en rempla\c{c}ant par la suite $S_\ZZ$ par le sch\'ema $S_\ZZ'$ introduit ci-dessous.

Pour toute fibre g\'eom\'etrique $Y_{\bar{t}}$ de $g$, on a $H^1(Y_{\bar{t}},\Oo)=H^2(Y_{\bar{t}},\Oo)=0$, de sorte que $\Pic_{\tau}(Y)$ est \'etale sur $T$. De plus, $\Pic(Y_{\bar{t}})\simeq\ZZ$, engendr\'e par la classe de l'unique faisceau inversible (tr\`es) ample $\Oo(1)$ de degr\'e $\prod a_i$. Localement sur $T$ pour la topologie \'etale, il existe donc sur $Y$ un faisceau inversible relativement ample de degr\'e $\prod a_i$, soit $\Oo(1)$. Localement sur $T$, il est unique. Puisque $H^i(Y_{\bar{t}},\Oo(k))=0$ (FAC), $g_*\Oo(k)$ est localement libre de formation compatible aux changements de base (\cite{9} p.~19 ou EGA~III); d\`es lors, $\Oo(1)$ est tr\`es ample et tout choix d'une base de $g_*\Oo(1)$ fournit un diagramme
\[
\xymatrix{
Y \ar@{^(->}[r]^{j} \ar[rd]_{g} & \PP^{n+\nu}_T \ar[d]^{p} \\
& T
}
\]
avec $j^*\Oo(1)=\Oo(1)$.

Prouvons que $j(Y)$ est une intersection de type (i); (iii) r\'esultera de l'unicit\'e de $\Oo(1)$. Les morphismes $H^0(\PP^{n+\nu},\Oo(k))\to H^0(Y_{\bar{t}},\Oo(k))$ sont surjectifs, donc aussi $p_*\Oo(k)\to g_*\Oo(k)$ et le noyau de cet \'epimorphisme de faisceaux localement libres est localement libre; l'assertion en r\'esulte ais\'ement.

Soit $S_\ZZ'$ le sch\'ema sur $\ZZ$ qui param\'etrise les $\nu$-uples d'hypersurfaces $H_i\subset\PP^{n+\nu}$, avec $\deg(H_i)=a_i$. Soit $S_\ZZ''$ le sous-sch\'ema ouvert de $S_\ZZ'$ qui param\'etrise les $\nu$-uples $(H_i)$ avec $\bigcap H_i$ lisse de dimension $n$. Le sch\'ema $S_\ZZ''$ est lisse sur $\Spec(\ZZ)$, car ouvert d'un produit d'espaces projectifs. L'application \'evidente de $S_\ZZ''$ dans $S_\ZZ$ est surjective par (i). Elle est lisse (\`a l'aide de (i), on v\'erifie ais\'ement un crit\`ere infinit\'esimal de lissit\'e). D\`es lors, $S_\ZZ$ est lisse sur $\Spec(\ZZ)$.

Posons $S=S_\CC$, $X=X_\CC$ et soit $f\colon X\to S$ la famille universelle \ref{1.8}.

\subsection{}

\begin{lemma}[Lemme fondamental]\label{lem:1.10}
Soit $a\colon A\to S$ un sch\'ema ab\'elien sur $S$.
\begin{enumerate}
\item[(i)] S'il existe un nombre premier $\ell$ et un isomorphisme de $\QQ_\ell$-faisceaux
\[
\nu_\ell\colon R^1u_*\QQ_\ell\xrightarrow{\sim}R^1a_*\QQ_\ell,
\]
ou s'il existe un isomorphisme horizontal
\[
\nu_{\mathrm{DR}}\colon\mathscr{H}^1_{\mathrm{DR}}(J(X/S)/S)\xrightarrow{\sim}\mathscr{H}^1_{\mathrm{DR}}(A/S),
\]
alors $A$ est isomorphe \`a $J(X/S)$. L'isomorphisme
\[
\nu'\colon J(X/S)\xrightarrow{\sim}A
\]
est unique au signe pr\`es.
\item[(ii)] Si de plus $A$ est muni d'une polarisation de la s\'erie principale, alors $\nu'$ est compatible aux polarisations. Si $\nu_\ell$ \textup(resp.~$\nu_{\mathrm{DR}}$\textup) comme en \textup{(i)} est compatible aux formes de polarisation, alors $\nu_\ell$ \textup(resp.~$\nu_{\mathrm{DR}}$\textup) est induit par $\nu'$ ou par $-\nu'$.
\end{enumerate}
\end{lemma}

Nous nous appuierons sur le th\'eor\`eme suivant de Lefschetz.

\begin{itheorem}[Lefschetz]\label{thm:1.11}
Soient $s\in S$ et $X_s=f^{-1}(s)$.
\begin{enumerate}
\item[(i)] La repr\'esentation de monodromie de $\pi_1(S,s)$ sur $H^n(X_s,\QQ)$ est absolument irr\'eductible.
\item[(ii)] Pour tout nombre premier $\ell$, la repr\'esentation de monodromie de $\pi_1^{\mathrm{\acute{e}t}}(S,s)$ sur $H^n(X_s,\ZZ/\ell)$ est absolument irr\'eductible.
\end{enumerate}
\end{itheorem}

Pour un expos\'e en cohomologie $\ell$-adique de ce th\'eor\`eme, suivant de tr\`es pr\`es la d\'emonstration de Lefschetz (cf.~\cite{6} Ch~V, III~p.~106), voir SGA~7~bis.

Prouvons \ref{lem:1.10}. Consid\'erons les repr\'esentations de monodromie
\begin{align*}
\rho_S&\colon\pi_1(S,s)\to\GL(H^1(J(X_s),\ZZ))\simeq\GL(H^n(X_s,\ZZ))\\
\rho_A&\colon\pi_1(S,s)\to\GL(H^1(A_s,\ZZ)).
\end{align*}
L'hypoth\`ese de (i) implique que, apr\`es extension des scalaires \`a $\QQ_\ell$ ou $\CC$, ces repr\'esentations deviennent isomorphes. L'isomorphie de deux repr\'esentations est invariante par extension du corps de base, d'o\`u l'existence d'un isomorphisme $\pi_1(S,s)$-\'equivariant
\[
\nu''\colon H^1(J(X_s),\QQ)\xrightarrow{\sim}H^1(A_s,\QQ).
\]
D'apr\`es \ref{thm:1.11}~(i), l'image de $\nu''^{-1}(H^1(A_s,\ZZ))\subset H^1(J(X_s),\QQ)$ dans $H^1(J(X_s),\ZZ/\ell)$ est tout ou rien. Il en r\'esulte qu'un multiple rationnel de $\nu''$ est un isomorphisme
\[
\nu'_\ZZ\colon H^1(J(X_s),\ZZ)\xrightarrow{\sim}H^1(A_s,\ZZ).
\]
D'apr\`es \ref{thm:1.11}~(i), $\nu'_\ZZ$ est uniquement d\'etermin\'e au signe pr\`es.

D'apr\`es \cite{3} 4.4.11 et 4.4.12, $\nu'_\ZZ$ provient d'un isomorphisme de sch\'ema ab\'elien
\[
\nu'\colon J(X/S)\xrightarrow{\sim}A,
\]
qui, comme $\nu'_\ZZ$, est unique au signe pr\`es. Rappelons le principe de la d\'emonstration. On interpr\`ete $\nu'_\ZZ$ comme une section du syst\`eme local $\Hom(R^1u_*\ZZ,R^1a_*\ZZ)$; d'apr\`es \ref{thm:1.11}~(i), c'est l'unique section de ce syst\`eme local, \`a un facteur pr\`es. On sait par ailleurs (\cite{3} 4.12) que le module de toutes les sections d\'efinit en chaque point de $s$ une sous-structure de Hodge de $\Hom(H^1(J(X_s),\ZZ),H^1(A_s,\ZZ))$. On en tire que $\nu'_\ZZ$ est en chaque point de $S$ de type de Hodge $(0,0)$, i.e. un morphisme de structures de Hodge, et l'assertion en r\'esulte.

Prouvons (ii). D'apr\`es \ref{thm:1.11}~(i), toute polarisation $\psi'$ de $J(X/S)$ est un multiple rationnel $b\cdot\psi$ de la polarisation donn\'ee. Pour $\psi'$ de la s\'erie principale, $b$ est un entier inversible positif, i.e. $\psi'=\psi$, ce qui prouve la premi\`ere assertion. Quant \`a la seconde, on a a priori par \ref{thm:1.11}~(i) $\nu_\ell=b\cdot\nu'_\ell$ avec $b\in\QQ_\ell^*$ (resp.~$\nu_{\mathrm{DR}}=b\cdot\nu'_{\mathrm{DR}}$ avec $b\in\CC^*$) et la compatibilit\'e aux formes $\psi_\ell$ ou $\psi_{\mathrm{DR}}$ fournit $b^2=1$, i.e.~$b=\pm 1$.

\subsection{}\label{1.12}

D'apr\`es \ref{thm:1.11}~(i), l'isomorphisme $\alpha_\ell$ est caract\'eris\'e au signe pr\`es par le fait qu'il transforme $\psi_\ell$ en $\varphi_\ell$ \eqref{eq:1.6.2}. Il serait facile d'adapter (en les simplifiant) les arguments de \cite{4} \S6 pour d\'eduire de $1/\ell$ que les intersections compl\`etes $V_n(\ua)$ sur les corps finis v\'erifient les conjectures de Weil. On prouverait de m\^eme le r\'esultat peut-\^etre plus g\'en\'eral suivant.

\begin{theorem}\label{thm:1.13}
Soit
\[
\xymatrix{
X \ar[rr]^{f} \ar[rd] & & \PP^1_V \ar[ld] \\
& \Spec(V) &
}
\]
une vari\'et\'e projective et lisse de dimension $2m+2$ sur un anneau de valuation discr\`ete $V$ d'in\'egale caract\'eristique \`a corps r\'esiduel fini $\FF_q$. Soient $X_0$ et $X_{\bar{\eta}}$ les fibres sp\'eciale et g\'en\'erique de $X$. Supposons que la ``nouvelle'' partie de la cohomologie des sections hyperplanes lisses $H$ de $X_{\bar{\eta}}$ soit de niveau $\le 1$ au sens suivant.
\begin{itemize}
\item[$(*)$] Pour $p+q=2m+1$ et $|q-p|>1$, le conoyau
\[
\Coker\bigl(H^q(X_{\bar{\eta}},\Omega^p)\to H^q(H,\Omega^p_H)\bigr)
\]
est nul.
\end{itemize}
Alors, pour toute section hyperplane lisse $H_0$ de $X_0$, les valeurs propres de l'endomorphisme de Frobenius de $\Coker\bigl(H^{2m+1}(X_0,\ZZ_\ell)\to H^{2m+1}(H_0,\ZZ_\ell)\bigr)$ sont des entiers alg\'ebriques dont tous les conjugu\'es complexes sont de valeur absolue $q^{\frac{2m+1}{2}}$. Ces entiers alg\'ebriques sont divisibles par $q^m$.
\end{theorem}

\section{Sur $\QQ$}

\subsection{}\label{2.1}

Soit $\ell_0$ un nombre premier. Pour toute extension $K$ de $\QQ$, consid\'erons le probl\`eme suivant:

\smallskip

$(J,K)$. Construire un sch\'ema ab\'elien polaris\'e de la s\'erie principale $u\colon J_K\to S_K$ sur $S_K$ (\ref{1.8}), muni d'un isomorphisme
\[
\alpha_{\ZZ_{\ell_0}}\colon R^nf_*\ZZ_{\ell_0}(m)\xrightarrow{\sim}R^1u_*\ZZ_{\ell_0}
\]
qui transforme $\psi_{\ell_0}$, d\'efini comme en \eqref{eq:1.6.2}, en la forme de polarisation.

\smallskip

Pour $K=\CC$, le probl\`eme $(J,K)$ a d'apr\`es \ref{lem:1.10} une solution unique \`a isomorphisme unique pr\`es.

\subsection{}\label{2.2}

Un principe g\'en\'eral veut qu'un probl\`eme $P$ d\'efini sur un sous-corps $k$ de $\CC$ et ayant sur $\CC$ une solution unique \`a isomorphisme unique pr\`es ait une solution (tout aussi unique) sur $k$. Ici, on trouve qu'il existe un sch\'ema ab\'elien $u\colon J_\QQ\to S_\QQ$ sur $S_\QQ$, muni d'une polarisation principale $\psi$ et d'un isomorphisme $\alpha_{\ell_0}$ comme en $(J,K)$. De plus, le triple $(J_\QQ,\psi,\alpha_{\ell_0})$ est unique \`a isomorphisme unique pr\`es.

\subsection{}\label{2.3}

Le principe \'enonc\'e plus haut se justifie ainsi.

a) Si le probl\`eme $P$ est de trouver un nombre $a$, $a$ ayant quelque propri\'et\'e g\'eom\'etrique, le principe est \'evident: $k$ est le corps des invariants de $\Aut(\CC/k)$, et une solution $a$ de $P$ sur $\CC$, \'etant unique, est invariante par $\Aut(\CC/k)$. ``G\'eom\'etrique'' signifie: ``invariant par extension des scalaires''.

b) Si le probl\`eme $P$ est de trouver un sous-espace vectoriel d'un espace vectoriel d\'efini sur $k$, ayant quelque propri\'et\'e g\'eom\'etrique, on applique de m\^eme Bourbaki, Alg.~Chap~2 \S8 Th~1 \`a $\Aut(\CC/k)$. Il n'est pas requis que $V$ soit de dimension finie.

c) Si le probl\`eme $P$ est de trouver un sous-sch\'ema $Y$ d'un sch\'ema affine $W$ d\'efini sur $k$, ayant quelques propri\'et\'es g\'eom\'etriques, on applique b): l'id\'eal de $Y$ est un sous-espace vectoriel de l'anneau de coordonn\'ees.

d) Le cas o\`u, dans c), $W$ ne serait pas affine se ram\`ene \`a c) par recollement.

e) Pour tout entier $b\ge 3$, soit $M_b$ le sch\'ema sur $S_\ZZ[1/b]$ dont les sections sur tout $S_\ZZ[1/b]$-Sch\'ema $t\colon T\to S_\ZZ[1/b]$ classifient les sch\'emas ab\'eliens polaris\'es de la s\'erie principale $a\colon A\to T$ munis d'un isomorphisme $t^*R^nf_*\ZZ/b(m)\xrightarrow{\sim}R^1a_*\ZZ/b$. (Cet isomorphisme \'etant compatible aux formes d'intersection.)

Une solution du probl\`eme $(J,K)$ s'interpr\`ete comme une section sur $S_K$ de $M_{\ell_0}=\varprojlim_c M_{c\ell_0}$. On applique d) \`a cette section, vue comme sous-sch\'ema de $(M_{\ell_0})_K$.

\subsection{}\label{2.4}

\begin{remark}
Lorsque le probl\`eme $P$ a une solution sur une extension finie de $k$, le principe \'enonc\'e plus haut se d\'emontre par descente galoisienne. On peut souvent se ramener \`a ce cas en utilisant le principe de l'extension finie EGA IV 9.1.1.
\end{remark}

Soit $\ell$ un nombre premier. Soient $K$ le corps des fonctions de $S_\QQ$ et $\bar{K}$ une cl\^oture alg\'ebrique de $K$. Tout $\ZZ_\ell$-faisceau $\mathscr{F}$ sur $S_\QQ$ d\'efinit une repr\'esentation $\ell$-adique de $\Gal(\bar{K}/K)$ sur la fibre g\'en\'erique g\'eom\'etrique $\mathscr{F}_{\bar{\eta}}$ de $\mathscr{F}$.

\subsection{}\label{2.5}

\begin{lemma}\label{lem:2.5}
Les repr\'esentations $\ell$-adiques de $\Gal(\bar{K}/K)$ d\'efinies par $R^nf_*\ZZ_\ell(m)$ et par $R^1u_*\ZZ_\ell$ ont le m\^eme caract\`ere.
\end{lemma}

Soit $U$ un ouvert de Zariski de $S_\ZZ$ o\`u $\ell$ et $\ell_0$ sont inversibles, contenant $S_\QQ$, et tel que $J_\QQ$ soit la restriction \`a $S_\QQ$ d'un sch\'ema ab\'elien $u'\colon J'\to U$. Les repr\'esentations $\ell$-adique \ref{2.5} de $\Gal(\bar{K}/K)$ sont alors non ramifi\'ees en dehors de $U$: elles se factorisent par une repr\'esentation de $\pi_1(U,\bar{\eta})$. De m\^eme pour leurs analogues $\ell_0$-adiques. Chaque point ferm\'e $s$ de $U$ d\'efinit une classe de conjugaison de substitutions de Frobenius $\varphi_s\in\pi_1(U,\bar{\eta})$. On a
\begin{equation}\label{eq:2.5.1}
\Tr(\varphi_s^{-1},(R^1u_*\ZZ_\ell)_{\bar{\eta}})=\Tr(\varphi_s^{-1},(R^1u_*\ZZ_{\ell_0})_{\bar{\eta}})\in\ZZ_\ell:
\end{equation}
c'est la trace de l'endomorphisme de Frobenius de $J'_s$. De m\^eme,
\begin{equation}\label{eq:2.5.2}
\Tr(\varphi_s^{-1},(R^nf_*\ZZ_\ell(m))_{\bar{\eta}})=\Tr(\varphi_s^{-1},(R^nf_*\ZZ_{\ell_0}(m))_{\bar{\eta}})\in\QQ
\end{equation}
les deux membres valent (d'apr\`es la formule des traces de Lefschetz en cohomologie $\ell$-adique)
\[
\frac{1}{q^m}\bigl(\#\PP^n(k(s))-\#X_s(k(s))\bigr).
\]
Par hypoth\`ese, on a
\[
\Tr(\varphi_s^{-1},(R^1u_*\ZZ_{\ell_0})_{\bar{\eta}})=\Tr(\varphi_s^{-1},(R^nf_*\ZZ_{\ell_0}(m))_{\bar{\eta}}),
\]
d'o\`u, par \eqref{eq:2.5.1} et \eqref{eq:2.5.2},
\[
\Tr(\varphi_s^{-1},(R^1u_*\ZZ_\ell)_{\bar{\eta}})=\Tr(\varphi_s^{-1},(R^nf_*\ZZ_\ell(m))_{\bar{\eta}}).
\]
On conclut en notant que la trace est une fonction continue sur $\pi_1(U,\bar{\eta})$ et que d'apr\`es le th\'eor\`eme de densit\'e de \v{C}ebotarev, les Frobenius $\varphi_s$ sont denses dans $\pi_1(U,\bar{\eta})$.

\subsection{}\label{2.6}

\begin{lemma}\label{lem:2.6}
Les repr\'esentations de $\Gal(\bar{K}/K)$ sur $R^nf_*\QQ_\ell(m)$ et sur $R^nf_*\ZZ/\ell(m)$ sont irr\'eductibles.
\end{lemma}

En effet, l'image de $\Gal(\bar{K}/K)$ dans $\GL(R^nf_*\QQ_\ell(m))$ ou $\GL(R^nf_*\ZZ_\ell(m))$ contient celle du groupe fondamental g\'eom\'etrique $\pi_1^{\text{\'et}}(S_\QQ,\bar{\eta})$. Il r\'esulte de \ref{thm:1.11} que l'action de ce dernier est d\'ej\`a irr\'eductible.

De \ref{lem:2.5} et \ref{lem:2.6} on tire que les repr\'esentations $\ell$-adiques de $\Gal(\bar{K}/K)$ sur $(R^nf_*\ZZ_\ell(m))_{\bar{\eta}}$ et sur $(R^1u_*\ZZ_\ell)_{\bar{\eta}}$ sont isomorphes. Tout isomorphisme entre ces repr\'esentations provient d'un isomorphisme
\[
\alpha'_\ell\colon R^nf_*\ZZ_\ell(m)\xrightarrow{\sim}R^1u_*\ZZ_\ell.
\]
D'apr\`es \ref{thm:1.11}, apr\`es extension des scalaires de $\QQ$ \`a $\CC$, $\alpha'_\ell$ devient un multiple $\ell$-adique de l'isomorphisme \eqref{eq:1.6.1}: $\alpha'_\ell=b\cdot\alpha_\ell$ avec $b\in\ZZ_\ell^*$. Sur $\QQ$, on pose $\alpha_\ell=b^{-1}\alpha'_\ell$.

\subsection{}\label{2.7}

Nous avons ainsi obtenu:
\begin{enumerate}
\item[a)] un sch\'ema ab\'elien polaris\'e de la s\'erie principale $J_\QQ$ sur $S_\QQ$;
\item[b)] pour tout nombre premier $\ell$, un isomorphisme compatible aux formes d'intersection
\[
\alpha_\ell\colon R^nf_*\ZZ_\ell(m)\xrightarrow{\sim}R^1u_*\ZZ_\ell;
\]
\item[c)] un isomorphisme $J_\QQ\times_{S_\QQ}S_\CC\xrightarrow{\sim}J(X_\CC/S_\CC)$ qui identifie les isomorphismes pr\'ec\'edents aux isomorphismes \eqref{eq:1.6.1}.
\end{enumerate}

\subsection{}\label{2.8}

Pour toute extension $K$ de $\QQ$, soit $\mathscr{W}_K$ le $K$-espace vectoriel des morphismes horizontaux
\[
\alpha'\colon\mathscr{H}^n_{\mathrm{DR}}(X_K/S_K)\to\mathscr{H}^1_{\mathrm{DR}}(J_K/S_K).
\]
On a $\mathscr{W}_K=\mathscr{W}_\QQ\otimes_\QQ K$. Pour $K=\CC$, $\mathscr{W}_\CC$ est de dimension un, et ses \'el\'ements non nuls sont des isomorphismes. On en d\'eduit qu'il existe un isomorphisme
\[
\alpha'\colon\mathscr{H}^n_{\mathrm{DR}}(X_\QQ/S_\QQ)\xrightarrow{\sim}\mathscr{H}^1_{\mathrm{DR}}(J_\QQ/S_\QQ).
\]
De plus, $\alpha'$ est unique \`a multiplication pr\`es par $2\in\QQ^*$. On v\'erifie sur $\CC$ que
\begin{enumerate}
\item[a)] $\alpha'$ respecte les filtrations de Hodge des deux membres: $\alpha'(F^{m+1})=F^1$;
\item[b)] $\alpha'$ transforme $\psi_{\mathrm{DR}}$ en un multiple de $\varphi_{\mathrm{DR}}$:
\[
\alpha'(\psi_{\mathrm{DR}})=\mu(\alpha')\cdot\varphi_{\mathrm{DR}}\qquad(\mu(\alpha')\in\QQ^*).
\]
\end{enumerate}
On a $\mu(2\alpha')=2^{-2}\mu(\alpha')$.

\subsection{}\label{2.9}

\begin{proposition}\label{prop:2.9}
Le multiplicateur $\mu(\alpha')$ est positif.
\end{proposition}

Ceci se v\'erifie apr\`es extension des scalaires de $\QQ$ \`a $\RR$.

\subsection{}\label{2.10}

\begin{conjecture}\label{conj:2.10}
Le multiplicateur $\mu(\alpha')$ est un carr\'e.
\end{conjecture}

Cette conjecture \'equivaut aux suivantes.
\begin{equation}\label{eq:2.10.1}
\text{Pour un choix convenable de $\alpha'$, on a $\alpha'(\psi_{\mathrm{DR}})=\varphi_{\mathrm{DR}}$.}
\end{equation}
\begin{equation}\label{eq:2.10.2}
\text{Sur $\CC$, l'isomorphisme d\'eduit de $\alpha'$ transforme cohomologie rationnelle en cohomologie rationnelle.}
\end{equation}
\begin{equation}\label{eq:2.10.3}
\text{Via 2.7 c), l'isomorphisme $\alpha_{\mathrm{DR}}$ (\ref{eq:1.7.1}) est d\'efini sur $\QQ$.}
\end{equation}
Que cette conjecture r\'esulte tant de la conjecture de Hodge que de la conjecture de Tate est une maigre consolation.

Rappelons que
\[
G=PGL_{n+\nu+1}
\]
agit sur $X_\QQ/S_\QQ$ (\ref{1.8}).

\subsection{}\label{2.11}

\begin{lemma}\label{lem:2.11}
Il existe une et une seule action de $G$ sur $J_\QQ/S_\QQ$, prolongeant l'action de $G$ sur $S_\QQ$. Pour cette action, les isomorphismes $\alpha_\ell$ (\ref{2.7} b)) sont \'equivariants.
\end{lemma}

Pour toute extension $K$ de $\QQ$, consid\'erons le probl\`eme suivant.

\smallskip

$(I,K)$. Soit $\mu^*\colon G_K\times J_K\to J_K$ l'action de $G_K$ sur $S_K$. Trouver un isomorphisme $a\colon\mu^*J_K\xrightarrow{\sim}\pr_2^*J_K$ qui d\'efinisse une action de $G_K$ sur $J_K/S_K$.

\smallskip

Puisque $G$ est connexe, les $\alpha_\ell$ sont automatiquement \'equivariants. De l\`a, ou de Mumford \cite{9} 6.1 p.~115 r\'esulte que $a$, s'il existe, est unique. D'apr\`es le principe \'enonc\'e en \ref{2.2}, il suffit donc de montrer que, sur $\CC$, $a$ existe. Ceci r\'esulte de ce que, sur $\CC$, la jacobienne interm\'ediaire est bien d\'efinie \`a isomorphisme unique pr\`es et de ce que $\mu^*(X_K/S_K)$ est donn\'e comme isomorphe \`a $\pr_2^*(X_K/S_K)$.

On d\'eduit de \ref{prop:1.9} et \ref{lem:2.11} le th\'eor\`eme suivant.

\subsection{}\label{2.12}

\begin{theorem}\label{thm:2.12}
\`A isomorphisme unique pr\`es, il existe un et un seul foncteur du type suivant $J$, muni des donn\'ees additionnelles suivantes.
\begin{enumerate}
\item[a)] Pour $T$ un sch\'ema de caract\'eristique $0$ et $g\colon Y\to T$ un morphisme propre et lisse, de fibr\'es des $V_n(\ua)$, $u\colon J(Y/T)\to T$ est un sch\'ema ab\'elien polaris\'e de la s\'erie principale.
\item[b)] $J$ est compatible aux changements de base $c\colon T'\to T$: des isomorphismes $J(c^*Y/T')\simeq c^*J(Y/T)$ sont donn\'es.
\item[c)] Des isomorphismes compatibles aux changements de base
\[
\alpha_\ell\colon R^ng_*\ZZ_\ell(m)\xrightarrow{\sim}R^1u_*\ZZ_\ell
\]
sont donn\'es.
\item[d)] Quand on se restreint aux sch\'emas sur $\CC$ lisses, le foncteur $(J,(\alpha_\ell))$ devient isomorphe au foncteur ``jacobienne interm\'ediaire'' construit au \S\,1.
\end{enumerate}
\end{theorem}

\section{Sur $\ZZ$}

La th\'eorie, sur $\ZZ$, est aussi peu satisfaisante qu'en cohomologie de De Rham (\ref{2.8} -- \ref{2.10}).

\subsection{}\label{3.1}

\begin{proposition}\label{prop:3.1}
Il existe un morphisme propre et birationnel $b\colon\tilde{S}\to S_\ZZ$ tel que
\begin{enumerate}
\item[a)] $b$ induit un isomorphisme de $S_\QQ$ avec $S_\QQ$;
\item[b)] il existe un sch\'ema ab\'elien polaris\'e de la s\'erie principale $\tilde{J}$ sur $\tilde{S}$, dont la restriction \`a $S_\QQ$ soit $J_\QQ$.
\end{enumerate}
\end{proposition}

On utilise le lemme suivant.

\subsection{}\label{3.2}

\begin{lemma}\label{lem:3.2}
Soient $V$ un anneau de valuation discr\`ete de corps de fraction $K$ de caract\'eristique $0$ et $g\colon\Spec(V)\to S_\ZZ$ un morphisme. Alors, la vari\'et\'e ab\'elienne sur $K$ image r\'eciproque de $J_\QQ$ a bonne r\'eduction.
\end{lemma}

Soient $\bar{K}$ une cl\^oture alg\'ebrique de $K$, $X_V$ l'image r\'eciproque de $X$ sur $V$, $J_K$ l'image r\'eciproque de $J$ sur $\Spec(K)$, $X_{\bar{\eta}}$ et $J_{\bar{\eta}}$ les images r\'eciproques $X_V$ et $J_K$ sur $\bar{K}$.

Soit $k$ un entier $\ge 3$ inversible sur $V$. Puisque $X_V$ est lisse sur $V$, il r\'esulte du th\'eor\`eme de sp\'ecialisation en cohomologie $\ell$-adique que le $\Gal(\bar{K}/K)$-module $H^n(X_{\bar{\eta}},\ZZ/k)(m)$ est non ramifi\'e. D'apr\`es \ref{2.7}, ce $\Gal(\bar{K}/K)$-module est isomorphe \`a $H^1(J_{\bar{\eta}},\ZZ/k)$, et il reste \`a appliquer le crit\`ere de bonne r\'eduction de N\'eron-Ogg-Chafarevitch.

Une m\'ethode br\`eve, mais trop compliqu\'ee, pour d\'eduire \ref{prop:3.1} de \ref{lem:3.2} est la suivante. Soit $k$ un entier $\ge 3$ et soit $M_k$ comme en \ref{2.3} e). Alors, $J_\QQ$ d\'efinit une section de $M_k$ sur $S_\QQ$. Soit $S_k$ l'adh\'erence de cette section dans $M_k$. Pour prouver que $b\colon S_k\to S_\ZZ[1/k]$ est propre, il suffit de tester le crit\`ere valuatif de propret\'e pour les anneaux de valuations discr\`etes de corps des fractions de caract\'eristique $0$: en effet, $S_{k\QQ}$ est dense dans $S_k$. Pour ceux-ci, il suffit d'appliquer \ref{lem:3.2}.

On v\'erifie ais\'ement que, au-dessus de $\Spec(\ZZ[1/k_1k_2])$, $S_{k_1}$ co\"incide avec $S_{k_2}$. Les $S_k$ se recollent en le $\tilde{S}$ cherch\'e.

\subsection{}\label{3.3}

Je conjecture que, dans \ref{prop:3.1}, on peut prendre $\tilde{S}=S_\ZZ$. Si tel \'etait le cas, \ref{thm:2.12} resterait vrai sur $\ZZ$.

On v\'erifie ais\'ement:

\subsection{}\label{3.4}

\begin{proposition}\label{prop:3.4}
Soit $a\colon\tilde{J}\to\tilde{S}$ comme en \ref{prop:3.1}. Pour chaque nombre premier $\ell$, l'isomorphisme $\alpha_\ell$ (\ref{2.7}) se prolonge sur $\tilde{S}[1/\ell]$ en
\[
\alpha_\ell\colon R^nf_*\ZZ_\ell(m)\xrightarrow{\sim}R^1a_*\ZZ_\ell.
\]
\end{proposition}

\subsection{La conjecture de Weil pour les $V_n(\ua)$}\label{3.5}

A l'aide de \ref{prop:3.4}, on peut donner une variante de la d\'emonstration indiqu\'ee en \ref{1.12} de la conjecture de Weil pour les $V_n(\ua)$.

Soit $X$ une intersection compl\`ete lisse de dimension $n$ et de multidegr\'e $\ua$ dans $\PP^{n+\nu}(\FF_q)$. $X$ correspond \`a un point $x\in S_\ZZ(\FF_q)$.
\begin{enumerate}
\item[a)] Il existe un point $\tilde{x}\in\tilde{S}(\FF_q)$ au-dessus de $x$.
\end{enumerate}
En effet, on peut relever $X$ en une intersection compl\`ete sur les vecteurs de Witt $W(\FF_q)$, ce qui rel\`eve $x$ en $x'\in S_\ZZ(W(\FF_q))$. Le point $x'$ d\'efinit un point de $S_\ZZ(W(\FF_q)\otimes\QQ)$ et, puisque $b$ est propre, ce dernier se prolonge en $\tilde{x}'\in\tilde{S}(W(\FF_q))$. On d\'eduit $\tilde{x}$ de $\tilde{x}'$ par r\'eduction.
\begin{enumerate}
\item[b)] Soit $\tilde{J}$ la vari\'et\'e ab\'elienne sur $\FF_q$, fibre de $\tilde{J}$ en $\tilde{x}$. Si $\bar{X}$ et $\bar{\tilde{J}}$ se d\'eduisent de $X$ et $\tilde{J}$ par extension des scalaires de $\FF_q$ \`a une de ses cl\^otures alg\'ebriques $\bar{\FF}_q$, \ref{prop:3.4} fournit un isomorphisme
\[
H^n(\bar{X},\ZZ_\ell)(m)\simeq H^1(\bar{\tilde{J}},\ZZ_\ell).
\]
\end{enumerate}
Cet isomorphisme commute aux endomorphismes de Frobenius (g\'eom\'etriques), de sorte que
\begin{equation}\label{eq:3.5.1}
\det(1-Ft,H^n(\bar{X},\ZZ_\ell))=\det(1-q^mFt,H^1(\bar{\tilde{J}},\ZZ_\ell)).
\end{equation}
La conjecture de Weil pour $X$ r\'esulte donc d'un th\'eor\`eme de Weil \cite{11} pour $\tilde{J}$.

\begin{thebibliography}{99}

\bibitem{1}
Clemens, C.H., Griffiths, P.A.: The intermediate jacobian of the cubic threefold ($\grave{\text{a}}$ para\^itre).

\bibitem{2}
Deligne, P.: Equations diff\'erentielles \`a points singuliers r\'eguliers. Lecture Notes in Mathematics 163. Berlin-Heidelberg-New York: Springer 1970. Je signale que II 1.23 et II 1.24 de [2] sont faux. La d\'emonstration de 4.1 doit \^etre modifi\'ee. J'envoie un erratum \`a quiconque me le demande.

\bibitem{3}
Deligne, P.: Th\'eorie de Hodge II. Publ. Math. IHES 40 (to appear).

\bibitem{4}
Deligne, P.: La conjecture de Weil pour les surfaces K3. \textit{Inventiones math.} \textbf{15}, 206--226 (1972).

\bibitem{5}
Griffiths, P.A.: Period of integrals on algebraic manifolds: Summary of main results and discussion of open problems. \textit{Bull. AMS} \textbf{76}, 228--296 (1970).

\bibitem{6}
Lefschetz, S.: L'analysis situs et la g\'eom\'etrie alg\'ebrique. Paris: Gauthier-Villars 1924. Reproduit dans Selected Papers. New York: Chelsea publ. Co. 1971.

\bibitem{7}
Lieberman, D.: Higher Picard varieties. \textit{Am. J. Math.} \textbf{90}, 1165--1199 (1968).

\bibitem{8}
Manin, Y.: Correspondences, motives and monoidal transformation. \textit{Matemati\v{c}eski Sbornik} \textbf{77} (119), 475--507.

\bibitem{9}
Mumford, D.: Geometric invariant theory. Ergebnisse 34. Berlin-Heidelberg-New York: Springer 1965.

\bibitem{10}
Rapoport, M.: Compl\'ement \`a l'article de P. Deligne: ``La conjecture de Weil pour les surfaces K3''. \textit{Inventiones math.} \textbf{15}, 227--236 (1972).

\bibitem{11}
Weil, A.: Vari\'et\'es ab\'eliennes et courbes alg\'ebriques. Paris: Hermann 1948.

\bibitem{EGA}
EGA. El\'ements de g\'eom\'etrie alg\'ebrique par A. Grothendieck et J. Dieudonn\'e. Publ. Math. IHES.

\bibitem{SGA7}
SGA 7 bis. S\'eminaire de g\'eom\'etrie alg\'ebrique 7-2e partie par P. Deligne et N. Katz. (Diffus\'e par l'IHES.)

\end{thebibliography}

\bigskip
\noindent
Pierre Deligne\newline
Institut des Hautes Etudes Scientifiques\newline
35, Route de Chartres\newline
F-91 Bures-sur-Yvette\newline
France

\medskip
\noindent
(Re\c{c}u le 26 ao\^ut 1971)

\end{document}
